\documentclass[11pt,a4paper]{article}
\usepackage[utf8]{inputenc}
\usepackage[T1]{fontenc}
\usepackage{lmodern}
\usepackage{amsmath,amssymb,amsfonts,amsthm}
\usepackage{geometry}
\geometry{margin=1in}
\usepackage{xcolor}
\usepackage{hyperref}
\usepackage{cleveref}
\usepackage{booktabs}
\usepackage{array}
\usepackage{microtype}
\usepackage{tcolorbox}

\hypersetup{
    colorlinks=true, linkcolor=blue, citecolor=red, urlcolor=cyan,
    pdftitle={The Selected K3 Surfaces of the Cooper s7/s10 Families: Identification, the K3xT2 Reading, and the Twisted-Weierstrass Route},
    pdfauthor={Xavier Callens and the SocrateAI Scientific Agora Collaboration}
}

\newtheorem{theorem}{Theorem}[section]
\theoremstyle{remark}
\newtheorem{remark}[theorem]{Remark}

\newcommand{\tierA}{\colorbox{green!15}{\textsf{\scriptsize Tier A}}}
\newcommand{\tierB}{\colorbox{yellow!20}{\textsf{\scriptsize Tier B --- derived, not measured}}}
\newcommand{\tierL}{\colorbox{blue!10}{\textsf{\scriptsize Tier L --- literature}}}
\newcommand{\tierC}{\colorbox{orange!15}{\textsf{\scriptsize Tier C --- blocked physics}}}
\newcommand{\lean}[1]{\texttt{#1}}

\newtcolorbox{headlinebox}{colback=red!4, colframe=red!45!black, boxrule=0.5pt, arc=1pt,
  left=6pt, right=6pt, top=4pt, bottom=4pt, fonttitle=\bfseries\small, title={Headline result}}
\newtcolorbox{scopebox}{colback=gray!5, colframe=gray!45, boxrule=0.4pt, arc=1pt,
  left=6pt, right=6pt, top=3pt, bottom=3pt, fonttitle=\bfseries\small, title={What is not claimed}}
\newtcolorbox{correctionbox}{colback=orange!6, colframe=orange!55!black, boxrule=0.5pt, arc=1pt,
  left=6pt, right=6pt, top=4pt, bottom=4pt, fonttitle=\bfseries\small, title={Corrections recorded in band}}

\title{\Large \textbf{The Selected $K3$ Surfaces of the Cooper $s_7$/$s_{10}$ Families:}\\[0.2em]
\textbf{Identification, the $K3\times T^2$ Reading, and the Twisted-Weierstrass Route}}
\author{\textbf{Xavier Callens} \and \textbf{SocrateAI Scientific Agora Collaboration} \\[0.3em]
\textit{Stream 2 --- Selection \& Geometry (\texttt{SocrateAI-Scientific-Agora-K3-DarkMatter})}}
\date{2026-10-08 \quad (preprint; not reviewed; continues the 2026-09-29 note)}

\begin{document}
\maketitle

\begin{scopebox}
This is a mathematics note. It makes no physical claim: no dark-matter identification, no dark-energy
statement, no observable. No point of either family maps to any measurement, the flux/tadpole route stays
blocked until a threefold base is specified, and the physical reading of ``$K3\times T^2$'' stays \tierC.
Nothing here ranks the $s_7$ family against the $s_{10}$ family. No Kodaira fibre type is attached to any
point of either family; we report discriminant-root orders and root lattices only.
\end{scopebox}

\begin{abstract}
The 2026-09-29 note adopted, by an explicit ruling, a within-family selector on the $\rho=20$ points of
the Cooper $s_7$ and $s_{10}$ Picard--Fuchs families: the point of minimal $|\operatorname{disc}T|$. This
note records what happened next. First, the selected points are identified as named surfaces. Through the
Shioda--Inose correspondence between singular $K3$ surfaces and even positive-definite binary forms, and
Takatsu's tabulation, the $s_7$ pick ($T=A_2$, $\det T=3$) is the unique singular $K3$ surface $X_3$ and the
$s_{10}$ pick ($\det T=4$) is the unique $X_4$. The pair is the one called ``Vinberg's two most algebraic $K3$
surfaces'' in the Sch\"utt--Shioda survey. Second, we state the one reading of ``$K3\times T^2$'' under which
a selection is meaningful, the Shioda--Inose pairing with a pair of $n$-isogenous elliptic curves, and
report the unscored consistency check between the level read from the lattice and the level read from the
modular side. Third, we record the twisted-Weierstrass route. A degree question that had conflated two
quantities is resolved, a two-divisor degree screen on candidate threefold bases is reported, and the
Mordell--Weil section required of the $M_7$-polarized member is computed. It has height $14$ and
$\bar P\cdot\bar O=5$ generically and at $z=1/27$, and $\bar P\cdot\bar O=0$ at $z=-1$. The selected point
$z=\infty$ has Mordell--Weil rank $0$ and carries no such section. At $z=1/27$ an explicit endomorphism of
square $-7$ is written over $\mathbb{Q}$. Every table is generated from a certificate.
\end{abstract}

\tableofcontents

\section{Scope}\label{sec:scope}

Ledger rules of this programme bound everything below. The $\rho=20$ cut and the discriminant-floor
selector are naming conventions inside one family, adopted by explicit rulings; they do not compare the
two families. The agreement between the binary-form side and the modular side of any statement here is
forced by the Shioda--Inose structure (an elliptic point of the modular curve is a CM point) and is never
offered as corroboration. The $s_{10}$ lattice certificate is LIVE since 2026-09-29; that changes no value
below and ranks nothing. The next step of the twisted-Weierstrass route that would involve a threefold
(work package TW3) has not been opened, because no ruling specifies twist data over a threefold base.

\section{The selection, in one table}\label{sec:recap}

\Cref{tab:rows} repeats the six $\rho=20$ points of the 2026-09-29 note, regenerated from the current
certificates. The lattice arithmetic of every row is kernel-checked in Lean~4 (\tierA); the identification
of the orthogonal complement with the family's transcendental lattice $T$ is \tierB\ (Dolgachev \S7, Doran
Thm.~5.13, \tierL). The selector picks, in each family, the row of minimal $|D|$: $z=\infty$ in both.

\begin{table}[h]
\centering
\begin{tabular}{lllrrll}
\toprule
Family & $z$ & $v$ & $v^2$ & $D$ & $T$ (reduced) & Tier / sources \\
\midrule
\lean{cooper\_s7} & $-1$ & $(2,-4,1)$ & $-2$ & $-7$ & $[2,1;1,4]$ & Tier A, 2 source(s) \\
\lean{cooper\_s7} & $1/27$ & $(1,-1,0)$ & $-2$ & $-28$ & $[2,0;0,14]$ & Tier A, 2 source(s) \\
\lean{cooper\_s7} & $\infty$ & $(14,-14,-5)$ & $-42$ & $-3$ & $[2,1;1,2]$ & Tier A, 2 source(s) \\
\lean{cooper\_s10}\,(ADVISORY) & $-1/4$ & $(2,-6,1)$ & $-4$ & $-20$ & $[4,2;2,6]$ & Tier A, 1 source(s) \\
\lean{cooper\_s10}\,(ADVISORY) & $1/16$ & $(1,-1,0)$ & $-2$ & $-40$ & $[2,0;0,20]$ & Tier A, 1 source(s) \\
\lean{cooper\_s10}\,(ADVISORY) & $\infty$ & $(10,-10,-3)$ & $-20$ & $-4$ & $[2,0;0,2]$ & Tier A, 2 source(s) \\
\bottomrule
\end{tabular}

\caption{The six $\rho=20$ points. Compared with the 2026-09-29 build, the $s_{10}$ rows no longer carry an
``(ADVISORY)'' tag; see \cref{sec:corrections}, item 1.}\label{tab:rows}
\end{table}

\section{Which surfaces were selected}\label{sec:identification}

A singular $K3$ surface is one of Picard number $20$. Shioda and Inose, in the form we read in Shioda's
survey and in Kumar--Kuwata (Thm.~8.1), give a bijection between singular $K3$ surfaces up to isomorphism
and even positive-definite binary lattices up to $\mathrm{SL}_2(\mathbb{Z})$, the surface going to its
transcendental lattice. A selected point is therefore named by its $T$. Takatsu's paper on the singular $K3$
surfaces of discriminant $3$, $4$ and $7$ states that each of these discriminants carries exactly one
surface, called $X_3$, $X_4$, $X_7$, and tabulates their models. Our checker parses that table and the
uniqueness sentence from the pinned text and, independently, enumerates the reduced forms of each
determinant.

\begin{table}[h]
\centering
\begin{tabular}{lllrrl}
\toprule
Family & $z$ & $T$ (Gram) & $\det T$ & classes & Surface \\
\midrule
\lean{cooper\_s7} & $\infty$$^{*}$ & $[2,1;1,2]$ & $3$ & $1$ & $X_{3}$ \\
\lean{cooper\_s7} & $-1$ & $[2,1;1,4]$ & $7$ & $1$ & $X_{7}$ \\
\lean{cooper\_s7} & $1/27$ & $[2,0;0,14]$ & $28$ & $2$ & not named by $\det T$ alone \\
\lean{cooper\_s10} & $\infty$$^{*}$ & $[2,0;0,2]$ & $4$ & $1$ & $X_{4}$ \\
\bottomrule
\end{tabular}

\caption{Identification of the $s_7$ points and the $s_{10}$ pick. $^{*}$: the point picked by the adopted
selector. ``classes'': the number of reduced even positive-definite forms of that determinant, computed.
At $z=1/27$ two classes share $\det T=28$, so the determinant alone names no surface there.}\label{tab:id}
\end{table}

\begin{headlinebox}
The $s_7$ pick is the singular $K3$ surface $X_3$ ($T=A_2$, $\det T=3$); the $s_{10}$ pick is $X_4$
($T=\langle2\rangle\oplus\langle2\rangle$, $\det T=4$). These are the two smallest determinants any even
positive-definite binary lattice attains. The $s_{10}$ family cannot reach $\det 3$, because $A_2$ does not
occur in it (a lattice fact recorded in an earlier certificate, not a preference). \tierB\ for the
identification of $T$; \tierL\ for the bijection and Takatsu's table.
\end{headlinebox}

\begin{remark}
The name ``Vinberg's two most algebraic $K3$ surfaces'' for the surfaces of discriminant $3$ and $4$ is
taken from the Sch\"utt--Shioda survey (\tierL). Vinberg's 1983 paper itself was not fetched or read. The
explicit $M_7$-polarized model of the 2026-09-29 note gives $E_\omega\times E_\omega$ at $z=\infty$, which
matches Takatsu's construction of $X_3$ from $E_3\times E_3$; this match is forced by the Shioda--Inose
structure and is reported as consistency, not as independent evidence.
\end{remark}

\section{What ``$K3\times T^2$'' can mean here}\label{sec:k3t2}

Two readings were put to the decision authority on 2026-09-16. Under the literal product reading, the
$T^2$ factor is a fixed torus unrelated to the $K3$, and every $K3$ surface serves equally well, so no
selection follows from it. Under the Shioda--Inose reading, ``$T^2$'' is the elliptic side of the pairing the
programme has certified: the order-three operator is the symmetric square of an order-two operator, and a
member of Picard number $19$ is paired with a product $E\times E'$ of two elliptic curves related by a cyclic
isogeny of degree $n$, with $T\cong U\oplus\langle 2n\rangle$. Selection then becomes a question about $n$,
which can be checked. The 2026-09-16 proposal adopted this reading and the check below presupposes it, but
the choice between the two readings has \emph{not} been ruled by the decision authority (see
\cref{sec:corrections}, item 6). The isogeny statement is standard; its primary source (Morrison, 1984) was
not fetched for this note.

The check adopted for this reading is unscored: it compares the level $n$ read from the lattice certificate
(Gram determinant and a serialized $U$-splitting witness, both recomputed) with the level read from the
modular side (an eta-quotient Hauptmodul for $\Gamma_0(n)^+$ or $\Gamma_0(n)^*$, with Newman's conditions
and an Atkin--Lehner check). \Cref{tab:t3} gives the result. Its power to fail rests on real negative
cases: four order-three entries of the same register fail the modular leg.

\begin{table}[h]
\centering
\begin{tabular}{llll}
\toprule
Family & $n$ from lattice & $n$ from modular side & Verdict \\
\midrule
\lean{cooper\_s7} & $7$ & $7$ & T3\_AGREE(n=7) \\
\lean{cooper\_s10} & $10$ & $10$ & T3\_AGREE(n=10) \\
\bottomrule
\end{tabular}

\caption{Lattice level against modular level. An unscored consistency gate, never a ranking.}\label{tab:t3}
\end{table}

The physical reading of $K3\times T^2$ as a compactification is \tierC\ and is not used. A separately
submitted document linking $K3\times T^2$ to Mathieu moonshine was checked claim by claim before this note
was written: of twelve integer-sequence identifiers it cited, seven were fabricated or attached to the wrong
object. Moonshine plays no part here; in any case the $K3$ elliptic genus is the same for every $K3$ surface,
so it cannot separate candidates.

\section{The twisted-Weierstrass route}\label{sec:tw}

A strict-pullback realization of the $s_7$ family as a Calabi--Yau threefold over the bases tried was
excluded earlier (a structural negative, \tierB). The route kept open is a twisted Weierstrass model whose
$K3$ fibres carry the $M_7$ polarization. Three steps have results.

\subsection{TW0: a degree question that conflated two quantities}

An earlier plan set a fundamental line-bundle degree $\ell=2$ ``via the symmetric-square theorem''.
Recomputing from the operator's Riemann scheme separates two numbers (\cref{tab:tw0}). The degree of the
Weierstrass model's fundamental line bundle is $\chi(\mathcal{O}_{K3})=2$, from Noether's formula with
$\deg\Delta=24$ (\tierA-level arithmetic, standard). The family-level Hodge bundle over the orbifold curve
$X_0(7)^+$ has degree $1/3$, and its square has degree $2/3$ (\tierB). The earlier formula reproduced $1$
from a Fuchs-relation identity. The rule now reads $\ell:=\chi(\mathcal{O}_{K3})=2$.

\begin{table}[h]
\centering
\begin{tabular}{ll}
\toprule
Quantity & Value \\
\midrule
orbifold Euler characteristic of $X_0(7)^+$ & $-2/3$ \\
$\deg\omega$ (family Hodge bundle, $\mathbb{Q}$-degree) & $1/3$ \\
$\deg\omega^{2}$ & $2/3$ \\
$\chi(\mathcal{O}_{K3})$ (elliptic-surface reading of $\ell$) & $2$ \\
$\deg\Delta$ of the Weierstrass model & $24$ \\
\bottomrule
\end{tabular}

\caption{The two degrees that the earlier plan had merged, from the $s_7$ operator's exponents.}\label{tab:tw0}
\end{table}

\subsection{TW1: a necessary screen on threefold bases}

The fibres of the generic member contain two $E_8$ root configurations. A twisted model over a threefold base
$B_3$ needs two divisors carrying the corresponding Tate vanishing orders ($v(f)\ge4$, $v(g)\ge5$,
$v(\Delta)=10$) inside the degree budget set by $-4K$, $-6K$ and $-12K$. \Cref{tab:tw1} lists the result.
On $\mathbb{P}^3$ the budget alone is met, but any two effective divisors meet in a curve, along which the
orders reach $(8,10,20)$ and the model is not minimal; the screen fails. On $\mathbb{P}^1\times\mathbb{P}^2$
and on the projective bundles $\mathbb{P}(\mathcal{O}\oplus\mathcal{O}(n))$ over $\mathbb{P}^2$, two disjoint
divisors exist and the screen passes. Passing is a necessary condition only: nothing here constructs a model.

\begin{table}[h]
\centering
\begin{tabular}{ll}
\toprule
Base $B_3$ & Two-divisor degree screen \\
\midrule
$\mathbb{P}^3$ & FAIL \\
$\mathbb{P}^1\times\mathbb{P}^2$ & PASS \\
$\mathbb{P}$($\mathcal{O}$ $\oplus$ $\mathcal{O}(n)$) over $\mathbb{P}^2$ & PASS (degree budget all $n\ge0$; realizability checked $n\le18$) \\
\bottomrule
\end{tabular}

\caption{Two-divisor degree screen. The certificates were written on 2026-07-29 and their status string
still reads ``DRAFT''; the screen was made live by ruling D20$'$ on 2026-10-07 (\cref{sec:corrections}).}\label{tab:tw1}
\end{table}

\subsection{TW2: the section that the $M_7$ polarization requires}

On a $K3$ surface with an elliptic fibration having two $E_8$ root configurations, the trivial lattice has
rank $18$ and discriminant $\pm1$. Reaching $|\operatorname{disc}\mathrm{NS}|=2n$ with Picard number $19$
then needs one Mordell--Weil section $P$ with height $h(P)=2n$. The height formula (Sch\"utt--Shioda
Thm.~11.5) gives $h(P)=2\chi+2\,\bar P\cdot\bar O-\sum\mathrm{contr}_v(P)$, so with $\chi=2$ and no
contributions, $\bar P\cdot\bar O=n-2$: for $s_7$, $h=14$ and $\bar P\cdot\bar O=5$. Kumar--Kuwata give an
independent route for $E\not\cong E'$: the section attached to an isogeny of degree $d$ has height $2d$,
and $d=7$ gives the same $14$.

At the three $\rho=20$ points the two curves become isomorphic and the count changes. Exact reduction of the
explicit model at each point gives the extra root lattice; the height formula then fixes the section
(\cref{tab:tw2}).

\begin{table}[h]
\centering
\begin{tabular}{lrrlrll}
\toprule
$s_7$ member & $|\mathrm{disc}\,T|$ & $\rho$ & extra roots & MW rank & $h(P)$ & $\bar P\cdot\bar O$ \\
\midrule
generic & $14$ & $19$ & --- & $1$ & $14$ & $5$ \\
$z=1/27$ & $28$ & $20$ & $A_{1}$ & $1$ & $14$ & $5$ \\
$z=-1$ & $7$ & $20$ & $A_{1}$ & $1$ & $7/2$ & $0$ \\
$z=\infty$ & $3$ & $20$ & $A_{2}$ & $0$ & --- & no section \\
\bottomrule
\end{tabular}

\caption{The section required by the $M_7$ polarization: generic member and the three $\rho=20$ points of
$s_7$. Root lattices from exact reduction of the explicit model; heights from Sch\"utt--Shioda
Thm.~11.5 (\tierB).}\label{tab:tw2}
\end{table}

\begin{headlinebox}
The section with $\bar P\cdot\bar O=5$ exists on the generic member and at $z=1/27$. At $z=-1$ the extra
$A_1$ absorbs part of the height and $\bar P\cdot\bar O=0$. At $z=\infty$, the point the selector picked, the
extra lattice is $A_2$, the Mordell--Weil rank is $0$ and there is no section at all. The selected surface
$X_3$ therefore does not carry the twisted-route section; the selection and the route are separate questions.
\end{headlinebox}

At $z=1/27$ the elliptic curve has $j=16581375$ and complex multiplication by the order of discriminant
$-28$. Using V\'elu's formulas over $\mathbb{Q}$ we wrote an explicit $7$-isogeny $\varphi$ from this curve to
a curve with the same $j$, and checked $\varphi\circ\varphi$ against multiplication by $7$ at $60$ rational
points, more than the degree bound of $49$. The $x$-coordinates fix $\varphi\circ\varphi$ only up to sign,
$\pm[7]$. The sign is $-7$ because the endomorphism algebra is the imaginary quadratic field of discriminant
$-28$, which has no element of square $+7$ (\tierB, through the CM certificate). The remaining step, carrying
the graph of $\varphi$ through the Kumar--Kuwata construction down to the $K3$ surface and reading
$\bar P\cdot\bar O=5$ on the pinned model, is specified but not done.

\section{What this note gives the dark-matter stream}\label{sec:stream3}

Nothing physical. The selected surfaces $X_3$ and $X_4$, the $K3\times T^2$ reading and the section
computation are mathematics about lattices and elliptic surfaces; no CM point and no selected surface maps to
any dark-matter observable, and no number here may be used as a prior, a parameter or a fit input. The
companion Stream~3 note (repository \texttt{DarkMatterK3-Home.github.io}, 2026-10-08) reports calibration
checks and a literature recount for a mixed fuzzy-dark-matter search window; it uses no result of this note.

\section{Corrections recorded in band}\label{sec:corrections}

\begin{correctionbox}
\begin{enumerate}
\item The table generator of the 2026-09-29 note attached ``(ADVISORY)'' to the $s_{10}$ rows by family
  name. The PDF rebuilt on 2026-10-07 therefore still printed three tags, while its own dated correction said
  the tables showed none. The generator now reads each row's \texttt{advisory\_family} flag (false for every
  row), a control fails if the tag is keyed on anything else, and that note carries a dated item 7.
\item The same generator typed the discriminant-root orders of its fibration table as literals, against the
  rule that numbers are computed. The printed values were right; they are now read from the certificate, and
  a tampered order changes the table (control).
\item On 2026-10-07 a record stated that $\varphi\circ\varphi=[-7]$ was proved. The computation fixed only
  $\pm[7]$; the sign comes from the CM discriminant. The record carries a dated correction.
\item The author of the discriminant $3,4,7$ paper was first taken from a search-result snippet, which
  named other authors; it was corrected from the paper's own header before any certificate was emitted.
\item The TW1 certificates still read ``DRAFT -- pending'' although ruling D20$'$ made the screen live. They
  are left unchanged, as provenance; their status is the one the ledger records.
\item \emph{Dated correction, 2026-10-08, after release \texttt{v0.3.17}.} The first release of this note
  called the Shioda--Inose reading of $K3\times T^2$ the one ``which this programme uses''. That reading was
  adopted by a proposal and is presupposed by the unscored T3 check, but the decision authority has not
  ruled between it and the product reading; the 2026-09-21 ``decision 1'' concerned the $\rho=20$ cut, not
  this choice. \Cref{sec:k3t2} now says so.
\item \emph{Dated correction, 2026-10-08.} The repository's regression list contained paper-lint lines
  (for the 2026-09-29 note since that date, and for this one), but the continuous-integration runner executes
  its own fixed list and never read them, so no paper had been linted in CI. The table check and the paper
  lint are now in the runner itself. The first release of this note was linted locally only.
\end{enumerate}
\end{correctionbox}

\section{Reproducibility}

Every table is generated from \texttt{data/certificates/} by \texttt{scripts/render\_paper\_tables.py},
whose \texttt{--check} mode fails if any fragment drifts from its certificate; its controls include tampered
inputs for every new table. The prose was linted by \texttt{scripts/check\_paper\_tier\_language.py}, which
applies the programme's tier-language rules to the whole body and blocks Kodaira fibre-type tokens; it ships
a self-test with planted violations. Decision records (D19$'$--D28$'$) and the identification brief:
\begin{itemize}\small
\item \nolinkurl{briefs/T0_DECISIONS_2026_10_07_STREAM2_DELEGATED.md}
\item \nolinkurl{briefs/SELECTED_K3_IDENTIFICATION_2026_10_08.md}
\end{itemize}

\section*{Provenance}
Generated-by: Claude (Opus~5.5), Stream~2 \textbar\ Verified-by: \texttt{render\_paper\_tables.py --check} and
its controls; each cited checker's own controls \textbar\ Reviewed-by: N (preprint; not reviewed).

\begin{thebibliography}{9}
\bibitem{Takatsu} T.~Takatsu, ``On the geometry of singular $K3$ surfaces with discriminant 3, 4 and 7,''
  \textit{arXiv:1903.03054}.
\bibitem{Shioda} T.~Shioda, ``$K3$ surfaces and sphere packings,'' \textit{J. Math. Soc. Japan} \textbf{60}
  (2008); MPIM preprint 2007-137.
\bibitem{KK} A.~Kumar and M.~Kuwata, ``Elliptic $K3$ surfaces associated with the product of two elliptic
  curves,'' \textit{arXiv:1409.2931}.
\bibitem{SS} M.~Sch\"utt and T.~Shioda, ``Elliptic surfaces,'' \textit{arXiv:0907.0298}.
\bibitem{Dolgachev} I.~Dolgachev, ``Mirror symmetry for lattice polarized $K3$ surfaces,''
  \textit{arXiv:alg-geom/9502005}.
\bibitem{Doran} C.~F.~Doran, ``Picard--Fuchs uniformization: modularity of the mirror map and
  mirror-moonshine,'' \textit{arXiv:math/9812162}.
\bibitem{CDLW} A.~Clingher, C.~F.~Doran, J.~Lewis, U.~Whitcher, ``Normal forms, $K3$ surface moduli, and
  modular parametrizations,'' \textit{arXiv:0712.1880}.
\bibitem{Prev} X.~Callens and the SocrateAI Scientific Agora Collaboration, ``Certified lattice data and a
  within-family selector for the Cooper $s_7$/$s_{10}$ $K3$ families,'' Stream~2 note, 2026-09-29
  (\texttt{papers/stream2\_selection\_geometry\_2026\_09\_29.pdf}).
\end{thebibliography}

\end{document}
